\chapter{Two quadratic families}
\label{chap:families}

The main work in the proof of Theorem~\ref{thm:main} is, in a sense, to show that~$a$
has~$\phi_n$ for all $n \ge 2$. The vehicle for this is a prime $p \equiv 3 \bmod 4$ such that
modulo~$p$ the sequence~$(c_n)$ becomes stationary at index~$n + 1$. Such a prime exists when
$|v_n| \equiv 3 \bmod 4$ or $|v_n| \equiv 6 \bmod 8$. (The character induced by the Jacobi symbol
$\Leg{\cdot}{|v_n|}$ or $\Leg{\cdot}{|v_n/2|}$ would state this in a more ``constructive'' way,
without relying on a factorization.) This is the reason for the conditions (C1)--(C3): they
provide the required congruences, and in all other cases the congruences fail. It is also why
the result covers only a set of relative density~$1/2$ in $\QQ \setminus (-2, \alpha)$.

In some cases one can get a bit further: when~$|v_n|$ (or~$u_n$, defined below), even though
congruent to~$2 \bmod 8$, systematically factors as a product of two positive numbers one of
which can be shown to be $\equiv 3 \bmod 4$ or $\equiv 6 \bmod 8$. Theorem~\ref{thm:families}
exploits this idea, which is due to Li~\cite{li2021}, who obtained essentially the special case
of integral~$a$.

We begin with additional criteria that provide~$\phi_n$ in certain cases. For $n \ge 1$ let
\[ u_n = r^2 w_n^2 - r w_n s^{2^{n-1}} + (r + s) s^{2^n - 1} \]
(Definition~\ref{def:wseq}); then $c_{n+1} - c_n + 1 = u_n/s^{2^n}$ in lowest terms, and in
particular~$u_n$ and~$s$ are coprime (Lemmas~\ref{lem:reduction} and~\ref{lem:wseq-res}). Since
\[ c_{n+2} + c_n = (c_{n+1} + c_n)(c_{n+1} - c_n + 1) \]
(Lemma~\ref{lem:orbit}(3)), modulo a prime dividing~$u_n$ we have $c_{n+2} \equiv -c_n$, and then
$c_{n+2k} \equiv -c_n$ for all $k \ge 1$ by induction.

\begin{lemma}[Period two]
  \label{lem:period2}
  \uses{lem:reduction, lem:orbit, lem:localeval, lem:wseq-res, lem:jump2}
  \lean{QuadraticIterates.hasCharPattern_mulSingle_of_dvd_uSeq,
    QuadraticIterates.cSeq_sq_sub_cSeq_add_eq_zero_of_uSeq}
  \leanok
  Let $n \ge 1$ and let $p \equiv 3 \bmod 4$ be a prime with $p \mid u_n$ and $p \nmid r + s$.
  Then the values of~$\chi_p$ on $(c_m)_{m \ge n}$ are $(x, y, -x, y, -x, \dots)$, with
  $x = \chi_p(c_n)$, $y = \chi_p(c_{n+1}) \in \{\pm 1\}$. Consequently, if~$a$ has~$\phi_0$
  and some character alternates from~$n$ on, then~$a$ has~$\phi_n$.
  \formalnote{For a character~$\psi$ modulo any power $N = p^k$, $k \ge 1$, with $\psi(-1) = -1$,
  $N \mid u_n$ and $p \nmid r + s$; the statement gives~$\phi_n$ from~$\phi_0$ and an alternating
  character directly (Lemma~\ref{lem:jump2}).}
\end{lemma}
\begin{proof}
  \leanok
  $u_n$ is prime to~$s$, so $\bar c_{n+2} = -\bar c_n$ in~$\ZZ/p\ZZ$ by the considerations above,
  and $(\bar c_m)_{m \ge n} = (\bar c_n, \bar c_{n+1}, -\bar c_n, \bar c_{n+1}, \dots)$
  (Lemma~\ref{lem:orbit}(2)). These entries are units: $\bar c_n = 0$ would give $p \mid r w_n$,
  so $p \mid r + s$, and $\bar c_{n+1} = 0$ would give $\bar c_n = 1$ and then $\bar a = -1$, so
  again $p \mid r + s$. Finally, $\chi_p(-1) = -1$ (Lemma~\ref{lem:localeval}), and
  Lemma~\ref{lem:jump2} gives the consequence.
\end{proof}

\begin{lemma}[Odd levels]
  \label{lem:odd}
  \uses{lem:period2, lem:pat-dyadic, lem:jump2, lem:reduction, lem:signs, lem:window, lem:prime3,
    lem:pat-rs, lem:mod4}
  \lean{QuadraticIterates.hasCharPattern_mulSingle_of_odd,
    QuadraticIterates.hasCharPattern_mulSingle_of_prime_dvd_uSeq,
    QuadraticIterates.exists_prime_dvd_uSeq_of_odd, QuadraticIterates.uSeq_emod_eight_of_odd,
    QuadraticIterates.uSeq_two_mul_add_one_modEq, QuadraticIterates.uSeq_pos,
    QuadraticIterates.cSeq_sq_sub_cSeq_add_pos_of_le_neg_two,
    QuadraticIterates.cSeq_sq_sub_cSeq_add_pos_of_neg_one_lt,
    QuadraticIterates.cSeq_sq_sub_cSeq_add_pos_of_odd}
  \leanok
  Let $a \equiv 1 \bmod 4$ with $a \le -2$. Then~$a$ has~$\phi_n$ for every odd $n \ge 3$.
  \formalnote{More generally for $a \equiv 1 \bmod 4$ with $a \le -2$, or $a > -1$, or $a < -1$ and
  $c_3 < 0 < c_4$; the condition $a \equiv 1 \bmod 4$ is stated as~$s$ odd and
  $r \equiv s \bmod 4$. \lname{hasCharPattern\_mulSingle\_of\_prime\_dvd\_uSeq} is
  Lemma~\ref{lem:period2} with the alternating character~$\val_2$.}
\end{lemma}
\begin{proof}
  \leanok
  Let $n \ge 3$ be odd. In~$\ZZ/8\ZZ$, $\bar c_n = 5$ for $\bar a = 1$ and $\bar c_n = 1$ for
  $\bar a = 5$, and $s^{2^n} \equiv 1$, so $u_n \equiv \bar c_n^2 - \bar c_n + \bar a + 1 \equiv 6
  \bmod 8$ (Lemma~\ref{lem:reduction}). Also $u_n > 0$: $u_n / s^{2^n} = c_n^2 - c_n + a + 1$,
  which is positive for $a \le -2$, as $c_n \ge -a \ge 2$ (Lemma~\ref{lem:signs}) gives
  $c_n(c_n - 1) + a + 1 \ge (a + 1)^2$ (and also for $a > 0$, for $-1 < a < 0$, as $c_n < 0$, and in
  the window, as $c_n \le c_3 < -1$, Lemma~\ref{lem:window}). So a prime $p \equiv 3 \bmod 4$
  divides~$u_n$ (Lemma~\ref{lem:prime3}). It does not divide~$r + s$: otherwise, in~$\ZZ/p\ZZ$,
  $\bar a = -1$ and $(\bar c_m)_{m \ge 1} = (1, 0, -1, 0, -1, \dots)$, so
  $\bar c_n^2 - \bar c_n + \bar a + 1 = 2 \ne 0$. Lemma~\ref{lem:period2}, with the alternating
  character~$\val_2$ (Lemma~\ref{lem:pat-dyadic}) and with~$\phi_0$ from Lemma~\ref{lem:mod4},
  gives~$\phi_n$.
\end{proof}

\begin{theorem}[Two quadratic families]
  \label{thm:families}
  \uses{lem:twoadic, lem:reflection, lem:period2, lem:odd, lem:jump2, prop:Z, prop:O,
    prop:T, lem:klein, lem:pat-dyadic, lem:strongcong, lem:prime3, prop:tfae}
  \lean{QuadraticIterates.nonempty_mulEquiv_of_eq_neg_mul_add_one,
    QuadraticIterates.exists_prime_dvd_vSeq_neg_mul_add, QuadraticIterates.neg_vSeq_neg_mul_add_eq,
    QuadraticIterates.nonempty_mulEquiv_of_eq_neg_sq_add_add_one,
    QuadraticIterates.exists_prime_dvd_uSeq_neg_sq_add_add, QuadraticIterates.uSeq_eq_mul_of_add_eq}
  \leanok
  Let $\ell = j/k$ with coprime positive integers~$j$ and~$k$.
  \begin{enumerate}
    \item If $(j, k) \equiv (2, 1)$ or $(3, 3) \bmod 4$ and $a = -\ell(\ell + 1)$ with
          $a \le -2$, then~$\rho_a$ is surjective.
    \item If $(j, k) \equiv (1, 1)$ or $(2, 3) \bmod 4$ and $a = -(\ell^2 + \ell + 1)$ with
          $a \le -2$, then~$\rho_a$ is surjective.
  \end{enumerate}
  \formalnote{Two theorems, with $j, k \in \ZZ$ as parameters (named~$u$, $v$ there) and the
  parameter written as $-(u/v)(u/v + 1)$, resp.\ $-((u/v)^2 + u/v + 1)$. In (1), $a \le -2$ is
  stated as $k < j$ (in these classes $\ell \ne 1$), in (2) as $k^2 \le j(j + k)$. The formal
  proof of (1) needs no bound on~$W$: $F_2 = jW + V > 0$ as $w_n < 0$, and $F_1 > 0$ as
  $F_1 F_2 = -v_n > 0$. In (2) it treats only the class $(1, 1)$, with~$k$ of either sign: the
  substitution $(j, k) \mapsto (j + k, -k)$, i.e.\ $\ell \mapsto -1 - \ell$, leaves~$a$, $r$
  and~$s$ unchanged and maps the class $(2, 3)$ to $(1, 1)$.}
\end{theorem}
\begin{proof}
  \leanok
  \uses{lem:wseq-res, lem:signs, def:wseq, lem:reduction}
  It suffices that~$a$ satisfies two of (Z), (O) and (T) and has~$\phi_n$ for all $n \ge 2$
  (Lemma~\ref{lem:klein}, Proposition~\ref{prop:tfae}).

  (1) Here $r = -j(j + k)$, $s = k^2$, $a < -2$, $\eps_a = -1$, and $v_2(r) = 1$, so (Z1) holds.
  A prime $p \equiv 3 \bmod 4$ divides~$r$ (one of the positive integers~$j + k$ and~$j$ is
  $\equiv 3 \bmod 4$), so (O4) or (T2) holds. By Lemma~\ref{lem:twoadic}(1), $a$
  has~$\phi_n$ (given that we know it has~$\phi_0$) when $n \ge 2$ and~$|v_n|$ is divisible by a
  prime $p \equiv 3 \bmod 4$. Let $n \ge 2$, $W = -w_n$ and $V = k^{2^n - 1}$. Then
  \[ -v_n = F_1 F_2 \qquad \text{with} \qquad F_1 = (j + k) W - V \quad \text{and} \quad
     F_2 = j W + V \]
  (expand, with $s = k^2$). Since $c_n \ge -a$ (Lemma~\ref{lem:signs}) and
  $c_n = r w_n / k^{2^n}$, $W \ge k^{2^n - 2}$, so $F_1 \ge j k^{2^n - 2} > 0$ and $F_2 > 0$. We
  have $s \equiv 1 \bmod 8$, so $w_{m+1} \equiv r w_m^2 + 1$, and as $r \equiv 2 \bmod 4$
  and~$w_m$ is odd, $w_m \equiv r + 1 \bmod 8$ for $m \ge 2$; thus $W \equiv j(j + k) - 1$ and
  $V \equiv k \bmod 8$. This gives $F_1 \equiv 6 \bmod 8$ when $(j, k) \equiv (2, 1) \bmod 4$
  and $F_2 \equiv 6 \bmod 8$ when $(j, k) \equiv (3, 3) \bmod 4$. In both cases a prime
  $p \equiv 3 \bmod 4$ divides $|v_n| = F_1 F_2$ (Lemma~\ref{lem:prime3}).

  (2) Here $r = -(j^2 + jk + k^2)$, $s = k^2$, $r + s = -j(j + k)$, and $a \equiv 1 \bmod 4$, so
  (O2) holds. $|r| \equiv 3 \bmod 4$ is not a square, so (Z1) holds. Lemma~\ref{lem:odd}
  gives~$\phi_n$ for odd $n \ge 3$, and~$\val_2$ alternates on~$c_n$ from $n = 1$
  (Lemma~\ref{lem:pat-dyadic}). Let $n \ge 2$ be even and $V = k^{2^n - 1}$. Then
  $u_n = G_n H_n$ with
  \[ G_n = r w_n - (j + k) V = k^{2^n}(c_n - \ell - 1) \quad \text{and} \quad
     H_n = r w_n + j V = k^{2^n}(c_n + \ell), \]
  both positive as $c_n \ge -a = \ell^2 + \ell + 1$. By Lemma~\ref{lem:strongcong},
  $w_n = (r + s) T$ with $T \equiv s^{2^{n-1} - 2} \bmod (r + s)$, so~$j$ and~$j + k$ divide~$w_n$.
  Set $K = H_n / j = -r(j + k) T + V$ when $(j, k) \equiv (1, 1) \bmod 4$, and
  $K = G_n / (j + k) = -r j T - V$ when $(j, k) \equiv (2, 3) \bmod 4$; in both cases the divisor
  is odd and $K > 0$. Then:
  \begin{itemize}
    \item $K \equiv 3 \bmod 4$. In~$\ZZ/4\ZZ$, $\bar a = 1$, so $\bar c_n = 2$ for even~$n$, and
          $r w_n = s^{2^{n-1}} \bar c_n = 2$ (Lemma~\ref{lem:reduction}); also $V \equiv k \bmod 4$.
          In the first case $jK = H_n \equiv 2 + jk$, so $K \equiv 2 + k \equiv 3$; in the second
          $(j + k)K = G_n \equiv 2 - (j + k)k$, so $K \equiv 2 - k \equiv 3 \bmod 4$.
    \item $K$ is prime to~$r + s$. $K$ is odd. Let~$q$ be an odd prime dividing
          $r + s = -j(j + k)$; then $q \nmid k$, $r \equiv -k^2$ and
          $r T \equiv -k^{2^n - 2} \bmod q$. In the first case $K \equiv (j + k)k^{2^n - 2} + V$,
          which is~$2V$ if $q \mid j$ and~$V$ if $q \mid j + k$; in the second case
          $K \equiv j k^{2^n - 2} - V$, which is~$-V$ if $q \mid j$ and~$-2V$ if $q \mid j + k$.
          The claim follows, as~$q$ does not divide~$2V$.
  \end{itemize}
  So a prime $p \equiv 3 \bmod 4$ divides~$K$, hence~$u_n$, and $p \nmid r + s$
  (Lemma~\ref{lem:prime3}). Lemma~\ref{lem:period2} with~$\val_2$ and~$\phi_0$ gives~$\phi_n$.
\end{proof}

Theorem~\ref{thm:families} generalizes the results of~\cite{li2021}. The first case in the
corollary below improves on~\cite{li2021}, Theorem~3.3 (which covers even~$m$); the second case
is~\cite{li2021}, Theorem~3.9.

\begin{corollary}[Li's families]
  \label{cor:li2021}
  \uses{thm:families}
  \lean{QuadraticIterates.nonempty_mulEquiv_of_eq_neg_mul_add_one_of_emod_four_eq_two,
    QuadraticIterates.nonempty_mulEquiv_of_eq_neg_sq_add_add_one_of_emod_four_eq_one}
  \leanok
  If $m \in \ZZ_{\ge 0}$ and $a = -(4m + 2)(4m + 3)$ or $a = -((4m + 1)(4m + 2) + 1)$,
  then~$\rho_a$ is surjective.
  \formalnote{Stated, as Theorem~\ref{thm:families} with $k = 1$, for $a = -\ell(\ell + 1)$
  with $\ell \equiv 2 \bmod 4$ (\lname{nonempty\_mulEquiv\_of\_eq\_neg\_mul\_add\_one\_of\_emod\_four\_eq\_two})
  and $a = -(\ell^2 + \ell + 1)$ with $\ell \equiv 1 \bmod 4$
  (\lname{nonempty\_mulEquiv\_of\_eq\_neg\_sq\_add\_add\_one\_of\_emod\_four\_eq\_one}), for a positive
  integer~$\ell$, that is, $\ell = 4m + 2$ resp.\ $\ell = 4m + 1$.}
\end{corollary}
\begin{proof}
  \leanok
  Apply Theorem~\ref{thm:families} with $k = 1$: $j = 4m + 2$ in the class $(2, 1)$, resp.\
  $j = 4m + 1$ in the class $(1, 1)$, where $a \le -3$.
\end{proof}
